% source: blackman_ch04_dynamic_targets_start__ocr_latex_draft.tex
% converted_by: Codex
% date: 2026-06-01
% compile_status: passed (xelatex)
% known_issues: Виявлено багато попереджень від xelatex про відсутні українські glyph у шрифтах EC; для повної локалізації бажано підключити UA-сумісні шрифти. Формули залишено з OCR-маркерами перевірки.
\documentclass[11pt]{article}
\usepackage[a4paper,margin=1in]{geometry}
\usepackage[T1]{fontenc}
\usepackage[utf8]{inputenc}
\usepackage{amsmath,amssymb}
\usepackage{hyperref}
\setlength{\parskip}{0.5em}
\setlength{\parindent}{0pt}
\begin{document}
\section*{Проєкт OCR/LaTeX: blackman\_ch04\_dynamic\_targets\_start}
\textbf{Джерело:} Blackman\_Popoli\_Design\_and\_Analysis\_of\_Modern\_Tracking\_Systems\\
\textbf{Сторінки PDF:} 229--237\\
\textbf{Тема:} системи відстеження; моделі динаміки цілі; моделі білого шуму (white-noise); початок розділу реалізації Калмана\\
\textbf{Увага:} Каркас сформований за результатами OCR. Перед публікацією перевірте рівняння, таблиці й рисунки на зчитування з оригінального PDF.\\
\clearpage\subsection*{Джерельна сторінка PDF 229 \quad Середня впевненість OCR: 96.12}
Моделювання і відстеження
динамічних цілей
4.1 Вступ
У цьому розділі розглядаються деякі специфічні моделі динаміки, які
використовують для побудови фільтрів Калмана для маневрених цілей. Також
обговорюються способи проектування адаптивних фільтрів для маневрування.
Ці питання отримали
значну увагу, і існують сотні робіт, присвячених методам \_\_ відстеження маневруючих цілей.
Майже всі ці роботи заявляють про кращі результати, ніж будь-які інші методи,
й подають дані для підтвердження таких тверджень. Отже, на сьогодні
відсутній універсальний підхід, який би однаково найкраще працював для всіх задач.
Мета цього розділу — дати огляд корисних моделей динаміки цілі
та адаптивних фільтрів, що використовують ці моделі. Остаточні висновки
не робляться, але декілька методів детальніше розглянуто разом із порівняльними
результатами. Вибір методів для детального розгляду базується на попередньому
дослідженні з акцентом на застосування в радіолокаційному контрольно-перехоплювальному
режимі для високоманеврових літаків. Проте тут же наводяться загальні міркування
та посилання, щоби зацікавлений читач міг дослідити альтернативні підходи
і інші області застосування.
Розділи 4.2 і 4.3 описують моделі динаміки цілі та те,
як ці моделі включаються у фільтр Калмана. Ці моделі є будівельними блоками
для різних адаптивних схем фільтрації, що розглядаються далі в розділі.
Розділ 4.4 подає огляд підходів, які використовуються для адаптивної
фільтрації маневрування. Порівняльні дослідження продемонстрували
199
\clearpage\subsection*{Джерельна сторінка PDF 230 \quad Середня впевненість OCR: 93.86}
200 Design and Analysis of Modern Tracking Systems
значну ефективність методу IMM (інтеракційної багатомодельної фільтрації) для
різноманітних застосувань. Саме тому в розділі 4.5 подано докладну реалізацію IMM.
Розділ 4.6 подає порівняльне дослідження кількох адаптивних фільтрів і
демонструє покращення якості, зазвичай пов'язане з підходом IMM.
Зручно ілюструвати методи та подавати порівняльні результати для
спрощеної задачі відстеження у горизонтальній площині. Безпосереднє поширення
на 3D-відстеження подано в розділі 4.7.
Більшість літератури з відстеження висвітлює проблеми, пов'язані з маневруванням
літаків. Однак також важливу увагу приділено задачі відстеження тактичних
балістичних ракет (TBM), зокрема SCUD, і це потребує інших моделей динаміки
та дещо іншого підходу за використання IMM-фільтрів. Отже, розділ 4.8 описує
динамічні моделі й подає підхід IMM для відстеження TBM-цілей. Розділ 4.9 дає
підсумок і кілька висновків авторського досвіду щодо методів цього розділу.
4.2 Моделі динаміки цілі
Невизначеність у оцінюванні стану через випадкову динаміку цілі або помилкове
моделювання цієї динаміки зазвичай задається коваріаційною матрицею шуму процесу
$Q$. Вибір $Q$ зазвичай ґрунтується не менш на експерименті (наприклад через
Монте-Карло моделювання), ніж на точному фізичному моделюванні.
У цьому розділі подано методи визначення $Q$, а детальніше обговорено в [1].
Розглянемо систему, яка використовує стан позиції цілі $x$, швидкості $v$, і, можливо,
прискорення $a$ як стани. Зазвичай прискорення цілі зберігається кілька
вибіркових інтервалів (7). Тому концепція корельованого в часі прискорення,
який моделюється як перший порядок Марковського процесу, має природне обґрунтування і
лише привела до моделі маневрування Сінгера [2, 3].
Модель Сінгера, ймовірно, є стандартною моделлю маневрування цілі,
тому її розглянуто першою. Однак, як зазначено в [3] і показали результати
Фіцджеральда [4], часто не можливо отримати точну оцінку прискорення,
якщо вимірюється лише позиція (а не швидкість). Тому також розглядаються моделі
білого шуму прискорення. Нарешті, подаються деталізованіші моделі, що враховують
зв'язок прискорення цілі з усіма вимірами положення. Також описано,
як ці моделі реалізуються за допомогою EKF.
4.2.1 Модель прискорення Сінгера
Модель маневрування цілі Сінгера [2] приймає стан цілі як позицію,
швидкість і прискорення та припускає, що прискорення є першим порядком
\clearpage\subsection*{Джерельна сторінка PDF 231 \quad Середня впевненість OCR: 70.37}
`Modeling and Tracking Dynamic Targets 201
Марковського процесу виду
% CHECK_FORMULA: OCR-рядок для (4.1) виглядає псевдовичерпним; перевірити в PDF.
a(k+1) = pma(k) + /1-- p2, omr(h) (4.1)
де
% CHECK_FORMULA: OCR-рядок для означення $P_m$ виглядає спотвореним; перевірити в PDF.
Pm=e Pl, B=l/tm
де $T_m$ — стала часу маневрування, $r(\&)$ — випадкова величина з нормальним
розподілом із нульовим середнім та одиничною дисперсією.
Інтервал дискретизації: $T$.
% CHECK_FORMULA: OCR-формула переходу фільтра Сінгера потребує звірки з оригіналом.
Отже, матриця переходу фільтра Сінгера задається:
T p(-1+BT+ pn)
31 -- Pm) (42)
]
@?=;0 1
0 0 Pm
% CHECK_FORMULA: Недійсні символи в наближенні для $T\ll T_m$; перевірити за першоджерелом.
Типово інтервал дискретизації менший за сталу часу маневрування ($T_m < T$).
\textasciitilde{} so-that a second-order approximation frequently used is -
% CHECK_FORMULA: Експоненціальне наближення (ряд?) у (4.2) потребує ручного відновлення.
. E
1 7 t
=\_ T
?/0 1 TI-x)
0 0 Pm
% CHECK_FORMULA: Формула (4.3) для $Q$ некоректно розпізнана OCR-рушницями, перевірити.
Точний розв’язок для коваріаційної матриці збурення процесу (або
маневру) записано як [2]:
% CHECK_FORMULA: Запис матриці $Q$ у (4.3) некоректний за OCR; потребує відновлення.
Q= = 921 922-923 (4.3)
'' 14931 932 933\}
де
] \textasciitilde{}--2pT 2p? T? 272. --BT
7 = 5531 --? +2BT+ --2B?T? --4BTe
% CHECK_FORMULA: Вираз для $Q$ у (4.3b) містить спотворені символи, перевірити.
qu = agile PT + 1 --2?787 4+ 26 Te P? \_ 2BT + pT]
% CHECK_FORMULA: Формула (4.3c) нечітка; верифікувати за першоджерелом.
113 >= 7p! - e PT \_ 9B TeF?)
\clearpage\subsection*{Джерельна сторінка PDF 232 \quad Середня впевненість OCR: 80.74}
202 Design and Analysis of Modern Tracking Systems
% CHECK_FORMULA: Перелік симетричних елементів (4.3) зчитано некоректно.
-- 
922 = aps te P* -- 3 -- 6 7PT + 287]
1 \_
gs = saz le PT F120 PT)
1
= 1 fy, -- --2AT
433 ap! er' |
Існують кілька цікавих граничних випадків. Перший: для малих інтервалів
вибірки ($T\ll T_m$), (4.3) переходить до
; 202
lim Q=--*
pI-0 Tm
3
ai oft sh
Si of ool
Ny Nf ofa
Аналогічно, у цьому випадку матриця переходу (4.2) спрощується до ньютонівської
матриці:
0 |
Якщо стала часу маневрування суттєво менша за інтервал дискретизації
($t_m < T$), точна оцінка прискорення неможлива, тому фільтр
повинен використовувати лише позицію і швидкість як стани. У цьому разі перехідна
матриця та коваріація шуму процесу спрощуються до
% CHECK_FORMULA: Формула (4.4) для білого шуму потребує відновлення символів.
BT>0
(ee ee)
iT rm 7
-- -- 2 3 2
% CHECK_FORMULA: Запис (4.4) переривчастий; перевірити структуру матриці.
=|) I]. Q= 202m BoP (44)
; 2
Практичний досвід та результати [3, 4] свідчать, що використання стану
прискорення часто доцільне лише тоді, коли доступне вимірювання швидкості.
При вимірюванні лише позиції точну оцінку прискорення можна отримати лише
коли $\tau_a > 10^{-7}$, і в цьому випадку досить точні прогнози позиції і
швидкості часто досягаються без прямої оцінки прискорення. Тому для зниження
обчислювальних витрат, у разі відсутності вимірювань швидкості, варто
використовувати двостановий (позиція і швидкість) фільтр із перехідною матрицею
і формою коваріації процесного шуму, заданими в (4.4).
\clearpage\subsection*{Джерельна сторінка PDF 233 \quad Середня впевненість OCR: 91.11}
Modeling and Tracking Dynamic Targets | 203
Модель білого шуму розглянуто детальніше в наступному розділі; порівняння
між моделлю Сінгера та моделлю білого шуму наведено в розділі 4.6.
4.2.2 Моделі сталої швидкості та сталого прискорення з білим шумом
У цих моделях випадкові динаміки цілі розглядаються як білий шум, який
входить у систему:
% CHECK_FORMULA: Формула (4.5) для $\dot a(t)$ втратила форматування; перевірити.
a,(t) = w(t) (4.5)
де процес білого шуму $w(t)$ визначається:
Elw()]=0, E[w(4)w(t)] = qd 6(?--- 7)
Як показано в Бар-Шалому та Лі [1], коваріаційна матриця шуму процесу
записується як
a=[F aC (4.6a)
os
Очікувано, (4.6a) збігається з граничним випадком моделі Сінгера,
коли інтервал дискретизації $T$ значно більший за час кореляції маневрування ($T_m$).
У такому разі $q$ можна інтерпретувати як
q= 202m
Вибір $q$ можна розглядати як налаштування параметра, найкраще
проводити шляхом моделювання, як це робиться в аналізі коваріацій у розділі 7
[3]. Там показано, що вибір
q= 1002
мінімізує похибки відстеження для білого шуму прискорення при треканні цілі
з корельованим прискоренням. Результати демонструють практично відсутність різниці
в роботі між моделлю білого шуму прискорення і моделлю Сінгера, навіть якщо модель
Сінгера була узгоджена з реальними динаміками цілі.
Умови, за яких двостатний фільтр може замінювати трьохстатну модель,
детальніше обговорено в розділі 13.
\clearpage\subsection*{Джерельна сторінка PDF 234 \quad Середня впевненість OCR: 87.27}
204 - Design and Analysis of Modern Tracking Systems
Модель білого шуму можна поширити на модель сталого прискорення (CA),
яка іноді називається моделлю білого деривативного (white jerk) шуму, так що
% CHECK_FORMULA: Формула (4.6b) OCR-спотворена, перевірити.
a,(t) = w(t)
Коваріаційна матриця процесу має ту саму форму, що й у моделі Сінгера
за умови $T_m \gg T$:
T> 74 TF3
20 8 6 .
Q=|F LF Lig (4.6b)
73 T?
oc a fF
Другий клас моделей білого шуму визначається так, що
ціль виконує прискорення або приріст прискорення, які є
сталими протягом інтервалу дискретизації та незалежними між сусідніми
інтервалами [1]. Ці моделі мають парні степені $\Delta t$ на
діагоналі матриці коваріації шуму процесу $Q$, як наведено далі:
Модель білого шуму зі сталою прискореною складовою
=(74 7] , 47
де, за [1], наведене емпіричне співвідношення між $a$ і максимальною величиною
прискорення,
@,,, 1S
05am < Oy XS am
Модель прискорення Вінера з поступальним білим шумом
T4/4 73/2 T?/2
% CHECK_FORMULA: Формули (4.7a) і (4.7b) викривлені OCR; перевірити.
Q=| 7/2 T? T |o; (4.7b)
T?/2 = T 1
де знову за [1] $A_{a,\max}$ позначає максимальну очікувану зміну прискорення
на інтервалі дискретизації $T$,
0.5Aa,, < dy < Aap,
Перша модель (4.7a) припускає, що вхідне прискорення некорельоване від скану
до скану, тоді як друга модель (4.7b) описує випадкову прогулянку
для прискорення.
\clearpage\subsection*{Джерельна сторінка PDF 235 \quad Середня впевненість OCR: 91.91}
Modeling and Tracking Dynamic Targets 205
Форма шуму процесу, подана у (4.6a) та (4.6b), переважає
форму (4.7a) і (4.7b), оскільки величина шуму, що входить у систему за
інтервал часу ($\Delta t > 0$), стає невідчутною до дискретизаційного
кроку. Наприклад, при використанні $Q$ з (4.6a), перехідна матриця
[7
і поширення Калмана (3.25e) дає співвідношення, яке безпосередньо
перевіряється прямою підстановкою
% CHECK_FORMULA: Формула розгалуження $Q(T+T')$ некоректна; перевірити.
Qi2T) = QT) + Q(T )o!
Це співвідношення показує, що одна й та сама кількість шуму процесу
надходить у систему на інтервалі $2T$ незалежно від того, чи крок дискретизації
$T$, чи $2T$.
Вибір параметрів процесного шуму [q in (4.6) і a in (4.7)]
як правило, є евристичною процедурою й виконується шляхом імітації
очікуваних траєкторій цілі зі фільтром відстеження й підбором параметрів
для мінімізації інтегрального критерію похибки. Альтернативно можна
застосувати аналіз коваріацій, як у [3] та розділі 13. Як уже згадувалося,
Бар-Шалом і Лі [1] дають рекомендації з вибору параметрів для (4.7).
-- 
4.2.3 Моделі координованого повороту
Розглянуті вище моделі в основному призначені для систем відстеження, де фільтри
майже некорельовані між собою, наприклад, незалежне відстеження за
$x$, $y$, $z$ координатами. Насправді типовий маневр цілі, наприклад,
поворот літака за координатами повороту, призводить до сильно корельованого
руху між напрямками. Саме тому цей розділ подає модель координованого
повернення, яка буде врахована в пізніших фільтрах Калмана.
% CHECK_FORMULA: Формула (4.8) визнана нечіткою OCR; перевірити.
Згідно з Maybeck et al. [5] і розділами 1–3 Bar-Shalom [6], модель
швидкості сталого повороту дає зміну вектора прискорення $a$
через кутову швидкість $w$, вектор швидкості $v$ і білий шум $w$:
% CHECK_FORMULA: Формула (4.8) та (4.9) потребують відновлення.
a= --w'v+w (4.8)
Кутова швидкість повороту $w$, що використовується в (4.8), задається
виразом векторного добутку: $o = v \times a_l$ (49)
y2
\clearpage\subsection*{Джерельна сторінка PDF 236 \quad Середня впевненість OCR: 74.68}
206 Design and Analysis of Modern Tracking Systems
% CHECK_FORMULA: Вхідні рівняння (4.8),(4.9) нечіткі — відновити з оригіналу.
Рівняння (4.8) і (4.9) дають базові динамічні зв’язки, що описують
літак під час координованого повороту. Розглянемо випадок, коли
прискорення перпендикулярне до вектора швидкості, а рух відбувається в
горизонтальній площині. Далі наведено розширення до 3D-моделей.
Компоненти вектора швидкості цілі задаються через модуль швидкості $s$
та кут курсу $4$:
Vy =scosh, vy =ssinh (4.10)
Далі розглянемо спеціальний випадок прискорення, перпендикулярного
до вектора швидкості. Тоді швидкість $s$ лишається сталою, а двічі
диференціюючи (4.10), отримуємо:
% CHECK_FORMULA: Формула (4.10) OCR-пошкоджена; потрібна перевірка.
ay = Wy, ay = --swsinh=--wv,
dy = --sw' cosh = --w'v,
--\_ 5 2 2
ay = Wy, ay = --W*d,
де w = / = кутова швидкість повороту цілі.
Рівняння прогнозування стану виведено далі. За умови сталої швидкості
повороту @ = 4 екстраполювані компоненти положення цілі мають вигляд:
% CHECK_FORMULA: Формули (4.11)-(4.12) потребують ручного відновлення.
X
x(k+ 1) = xb) +f veel) dt
te \textasciitilde{} . :
% CHECK_FORMULA: Формула (4.11) некоректно розпізнана.
x)=
x(k) +sT[SWceosh -- CWsin 4]
= x(k) + T[SWv,(k) -- CWv,(\&)] \textasciitilde{} (4.11)
% CHECK_FORMULA: Формула (4.12) потребує перевірки.
--y(R+1) = y(2) +s5T[SWsinh+ CWeos h]
% CHECK_FORMULA: Формула (4.12) продовження.
= y(k) + T[CWv,(\&) + SWv,(2)] (4.12)
\textasciitilde{} де
inwT 1--- T
sy ?sine cw A cos @ (4.13)
oT wT
Аналогічно для швидкості
% CHECK_FORMULA: Формула (4.14) потребує верифікації.
ve(k+1) = 5 cos(h + WT) = v?(A) cos wT -- v,(4) sinwT
% CHECK_FORMULA: Формула (4.14) OCR-виправлення потрібне.
vy(k+ 1) =ssin(a+ oT) = v,(k) sinwT + v,(k) coswT (4.14)
\clearpage\subsection*{Джерельна сторінка PDF 237 \quad Середня впевненість OCR: 94.74}
Modeling and Tracking Dynamic Targets ? an 207
Рівняння (4.11)–(4.14) задають базову динаміку
для цілі, що виконує кругове обертання у горизонтальній площині. Стан
прискорення позиції не потрібний, оскільки прискорення по суті визначається
швидкістю повороту (або кутовою швидкістю) $w$. Тому, як уже описано
в розділі 4.3, нелінійний фільтр для відстеження руху цілі у горизонтальній площині
можна задати п'ятьма станами $(x, v_x, y, v_y, \dot\psi)$. Далі,
другий, декупльований фільтр використовується для відстеження вертикального
($z$) руху. Нарешті, як також буде розглянуто пізніше, альтернативною
презентацією станів є використання позиції, швидкості, кута курсу та швидкості
зміни курсу як станів.
Відповідно до [6, 9, 10] та, як буде детальніше розглянуто нижче, фільтр Калмана
можна визначити через загальну залежність (4.8) із $w$, визначеним (4.9).
Це вводить нелінійність, оскільки $w$ оцінюється за допомогою оцінених станів.
Ця параметризація прямо розширюється до 3D-моделі маневру з
стандартними декартовими змінними стану, як визначено в [9] і далі розглядається.
Основні принципи моделі координованого повороту, описані раніше,
інтегровані й розвинені в інших моделях. Наприклад, як буде розглянуто
надалі, Алуані і Блейр [11] ввели обмеження, що прискорення майже
перпендикулярне до швидкості ($a-v\approx0$), у процес відстеження як
псевдовимірювання за методами розділу 3. Друге, Рекер і МакГіллім [12] розробили
техніку визначення маневру та трансформації у координатну систему, спеціально
сконструйовану для відстеження цілей на коловій траєкторії.
Асео та Ардільйо [13] подають модель другого порядку Маркова, яка поєднує
першопорядковий Марківський прискорення моделі Сінгера з (4.1) і модель
координованого повороту (4.8). Детальна реалізація для високочастотної системи
відстеження однієї цілі описана та оцінена в [13].
4.2.4 Інші підходи до моделювання динаміки цілі
У ідеалі динаміка цілі має задаватися у цільово-орієнтованій системі координат
з однією віссю вздовж вектора швидкості цілі. Теоретично цей підхід ефективний
як для літаків, що виконують координований поворот, де прискорення зазвичай
перпендикулярне до вектора швидкості, так і для певних класів ракет, де
прискорення в основному збігається з вектором швидкості. Боглер [14] розробив
повноцінний фільтр Калмана в таких системах швидкості, а Блейр і Паркер [15]
розглядають, як процесний шум можна сформувати у цих системах і
перетворити в декартову систему відстеження. Проте ці методи, як і
псевдовимірювальні підходи з [11], добре працюють лише за умови точного
оцінювання швидкості, щоб цільово-орієнтовані координати були коректно
визначені. Тому можлива дивергенція за умов помилкової швидкості.
\end{document}

